% preamble-exam-zh.tex — 共享 preamble
% 用于所有 exam-zh 模板
%
% 样式策略：
%   屏幕 → 语义彩色（蓝=关键想法，红=易错，绿=路线，灰=提示/教师备注）
%   打印 → 自动降级为黑白（通过 print mode 或 monochrome 选项）

\documentclass{exam-zh}

% ---------- 基础配置 ----------
\examsetup{
  page/size = a4paper,
  page/show-head = false,
  page/show-foot = false,
  sealline/show = false,
  font = times,
  math-font = xits,
}

\usepackage{tcolorbox}
\usepackage{needspace}
\tcbuselibrary{skins,breakable}

% ---------- 打印/屏幕颜色切换 ----------
% 用法：\documentclass[monochrome]{exam-zh} 即可切换为纯黑白
% 注意：不要使用 \if@xxx 放在 makeatother 外部；这里改成普通条件名。
\newif\ifeduprintcolor
\eduprintcolortrue

\makeatletter
\@ifclasswith{exam-zh}{monochrome}{\eduprintcolorfalse}{}
\makeatother

% ---------- 语义颜色定义 ----------
% 屏幕：有色彩区分；打印：降级为灰度
\ifeduprintcolor
  % --- 屏幕彩色模式 ---
  \definecolor{edu-blue}{RGB}{47, 82, 122}       % 关键想法
  \definecolor{edu-blue-bg}{RGB}{243, 246, 252}
  \definecolor{edu-red}{RGB}{180, 40, 40}         % 易错提醒
  \definecolor{edu-red-bg}{RGB}{255, 245, 240}
  \definecolor{edu-green}{RGB}{34, 100, 60}       % 路线图
  \definecolor{edu-green-bg}{RGB}{242, 250, 245}
  \definecolor{edu-orange}{RGB}{160, 90, 0}       % 训练目标
  \definecolor{edu-orange-bg}{RGB}{255, 248, 235}
  \definecolor{edu-gray}{RGB}{100, 100, 100}      % 提示/教师备注
  \definecolor{edu-gray-bg}{RGB}{248, 248, 248}
  \definecolor{edu-purple}{RGB}{90, 50, 120}      % 思考题
  \definecolor{edu-purple-bg}{RGB}{248, 244, 252}
\else
  % --- 打印黑白模式 ---
  \definecolor{edu-blue}{gray}{0.15}
  \definecolor{edu-blue-bg}{gray}{0.96}
  \definecolor{edu-red}{gray}{0.25}
  \definecolor{edu-red-bg}{gray}{0.97}
  \definecolor{edu-green}{gray}{0.20}
  \definecolor{edu-green-bg}{gray}{0.97}
  \definecolor{edu-orange}{gray}{0.25}
  \definecolor{edu-orange-bg}{gray}{0.98}
  \definecolor{edu-gray}{gray}{0.40}
  \definecolor{edu-gray-bg}{gray}{0.97}
  \definecolor{edu-purple}{gray}{0.30}
  \definecolor{edu-purple-bg}{gray}{0.98}
\fi

% ---------- 自定义教学环境 ----------
% 共性：enhanced + breakable（跨页不断裂）

% 原题卡片 — 强边框，突出显示
\newtcolorbox{problemcard}{
  enhanced, breakable,
  colback=edu-blue-bg,
  colframe=edu-blue,
  boxrule=1.2pt,
  arc=0pt,
  left=3mm, right=3mm, top=3mm, bottom=3mm,
}

% 解题路线图 — 左边线 + 浅底
\newtcolorbox{routebox}{
  enhanced, breakable,
  colback=edu-green-bg,
  colframe=edu-green!30,
  boxrule=0pt,
  borderline west={2.5pt}{0pt}{edu-green},
  left=3mm, right=3mm, top=3mm, bottom=3mm,
}

% 关键想法 — 蓝色左边线
\newtcolorbox{keyideabox}{
  enhanced, breakable,
  colback=edu-blue-bg,
  colframe=edu-blue!30,
  boxrule=0pt,
  borderline west={2.5pt}{0pt}{edu-blue},
  left=3mm, right=3mm, top=3mm, bottom=3mm,
}

% 易错提醒 — 红色左边线
\newtcolorbox{mistakebox}{
  enhanced, breakable,
  colback=edu-red-bg,
  colframe=edu-red!20,
  boxrule=0pt,
  borderline west={3pt}{0pt}{edu-red},
  left=3mm, right=3mm, top=2.5mm, bottom=2.5mm,
}

% 提示 — 灰色虚线左边线
\newtcolorbox{hintbox}{
  enhanced, breakable,
  colback=edu-gray-bg,
  colframe=edu-gray!20,
  boxrule=0pt,
  borderline west={2pt}{0pt}{edu-gray, dashed},
  left=3mm, right=3mm, top=2mm, bottom=2mm,
}

% 思考题 — 紫色虚线左边线
\newtcolorbox{thinkbox}{
  enhanced, breakable,
  colback=edu-purple-bg,
  colframe=edu-purple!20,
  boxrule=0pt,
  borderline west={2pt}{0pt}{edu-purple, dashed},
  left=2.5mm, right=2.5mm, top=2.5mm, bottom=2.5mm,
}

% 教师备注 — 灰色左边线，小字
\newtcolorbox{teachernote}{
  enhanced, breakable,
  colback=edu-gray-bg,
  colframe=edu-gray!15,
  boxrule=0pt,
  borderline west={2pt}{0pt}{edu-gray!60},
  left=2.5mm, right=2.5mm, top=2.5mm, bottom=2.5mm,
  fontupper=\small,
  colupper=black!60,
}

% 训练目标 — 橙色左边线，小字
\newtcolorbox{traininggoal}{
  enhanced, breakable,
  colback=edu-orange-bg,
  colframe=edu-orange!15,
  boxrule=0pt,
  borderline west={1.5pt}{0pt}{edu-orange},
  left=2mm, right=2mm, top=1mm, bottom=1mm,
  fontupper=\small,
}

% 答题区 — 无边框，纯留白
\newcommand{\answerarea}[1][40mm]{%
  \vspace{2mm}%
  \begin{tcolorbox}[
    enhanced, breakable,
    colback=white,
    colframe=white,
    boxrule=0pt,
    height=#1,
    valign=top,
    bottom=0mm,
    after skip=0mm,
  ]
  \end{tcolorbox}%
}

% ---------- 字体设置 ----------
% exam-zh (ctexbook) already sets CJK fonts via ctex.
% Only override if specific fonts are needed.
% macOS: Songti SC, STHeiti, STFangsong
% Linux: SimSun, SimHei, FangSong

% ---------- 页面设置 ----------
% exam-zh (ctexbook) already loads geometry.
% Override margins if needed:
% \geometry{top=16mm, bottom=16mm, left=14mm, right=14mm}

\examsetup{
  question/show-answer = true,
  fillin/show-answer = true,
  solution/show-solution = show-stay,
}

% ---------- 教师版解答样式 ----------
\newtcolorbox{solutionblock}{
  enhanced, breakable,
  colback=edu-blue-bg,
  colframe=edu-blue!25,
  boxrule=0pt,
  borderline west={2pt}{0pt}{edu-blue!50},
  arc=0pt,
  left=3mm, right=3mm, top=2mm, bottom=2mm,
  before skip=2mm,
  after skip=3mm,
  fontupper=\small,
}

% 解答步骤编号
\newcounter{solstep}
\newcommand{\solstep}[2]{%
  \stepcounter{solstep}%
  \par\medskip
  \noindent\hangindent=6mm
  {\color{edu-blue}\bfseries\textcircled{\scriptsize\thesolstep}\enspace #1}\enspace #2\par
}

\begin{document}

\section*{两圆位置关系与“$d, R, r$ 知二推一”专题练习（教师版）}

\section*{一、选择题（每题 4 分，共 8 分）}


\begin{question}[points=4]
  已知两圆的半径分别为 2 和 5。如果它们的圆心距 $d$ 是关于 $d$ 的方程 $|d - 5| = 2$ 的解，那么这两个圆的位置关系是\fillin。
  \begin{choices}
    \item 外离
    \item 相交
    \item 外切或内切
    \item 内含
  \end{choices}
  \paren[C]
  \begin{solutionblock}
    解方程 $|d - 5| = 2$ 可得 $d - 5 = 2 \implies d = 7$，或 $d - 5 = -2 \implies d = 3$。两圆半径为 2 和 5。和为 $5+2=7$，差为 $5-2=3$。当 $d=7$ 时，圆心距等于半径之和，两圆**外切**；当 $d=3$ 时，圆心距等于半径之差，两圆**内切**。因此两圆的位置关系是“外切或内切”（即相切）。故选 C。
  \end{solutionblock}
\end{question}



\begin{question}[points=4]
  已知 $\odot O_1$ 的半径为 3，$\odot O_2$ 的半径为 5。若它们的圆心距 $d$ 满足不等式组 $2 < d < 8$，则这两圆的位置关系是\fillin。
  \begin{choices}
    \item 外切
    \item 相交
    \item 内切
    \item 内含
  \end{choices}
  \paren[B]
  \begin{solutionblock}
    已知两圆半径为 3 和 5，半径之差为 $5-3=2$，半径之和为 $5+3=8$。因为圆心距 $d$ 满足 $2 < d < 8$，即满足区间 $R-r < d < R+r$，所以两圆的位置关系为**相交**。故选 B。
  \end{solutionblock}
\end{question}


\section*{二、填空题（每题 4 分，共 12 分）}


\begin{question}[points=4]
  已知两圆的半径分别为 3 和 7。若这两圆相切，则它们的圆心距 $d = $\fillin。
  \paren[4 或 10]
  \begin{solutionblock}
    因为两圆相切，必须分类讨论：(1) 外切时，圆心距 $d = R+r = 7+3=10$；(2) 内切时，圆心距 $d = R-r = 7-3=4$。所以 $d$ 为 4 或 10。
  \end{solutionblock}
\end{question}



\begin{question}[points=4]
  已知 $\odot A$ 的半径为 5，$\odot A$ 与 $\odot B$ 相切，两圆的圆心距 $d = 3$，则 $\odot B$ 的半径为\fillin。
  \paren[2 或 8]
  \begin{solutionblock}
    设 $\odot B$ 的半径为 $r$。因为相切，分类讨论：(1) 外切时，$d = r_A + r \implies 3 = 5 + r \implies r = -2$（因为半径必须为正数，故舍去）；(2) 内切时，由于谁是大圆未知，继续分类：①若 $\odot A$ 是大圆，$d = r_A - r \implies 3 = 5 - r \implies r = 2$；②若 $\odot B$ 是大圆，$d = r - r_A \implies 3 = r - 5 \implies r = 8$。因此，半径为 2 或 8。
  \end{solutionblock}
\end{question}



\begin{question}[points=4]
  已知两圆的半径分别为 3 和 5。若这两圆没有公共点，则两圆的圆心距 $d$ 的取值范围是\fillin。
  \paren[$d > 8$ 或 $0 \le d < 2$]
  \begin{solutionblock}
    两圆半径为 3 和 5。没有公共点说明它们“外离”或“内含”。外离时 $d > 5+3=8$；内含时 $0 \le d < 5-3=2$。故 $d$ 的范围是 $d > 8$ 或 $0 \le d < 2$。
  \end{solutionblock}
\end{question}


\section*{三、解答题（每题 10 分，共 30 分）}

\needspace{8\baselineskip}

\begin{problem}[points=10]
  在平面直角坐标系中，$\odot A$ 的圆心坐标为 $A(-3, 0)$，半径为 $1$；$\odot B$ 的圆心坐标为 $B(3, 0)$，半径为 $2$。
\begin{enumerate}[label=(\arabic*)]
  \item 求两圆初始位置时的圆心距 $d$ 及它们的位置关系；
  \item 若 $\odot A$ 沿 $x$ 轴正方向以每秒 $1$ 个单位长度的速度做平移运动，设运动时间为 $t$ 秒（$t > 0$）。当 $\odot A$ 与 $\odot B$ 相切时，求 $t$ 的值。
\end{enumerate}

  \answerarea[70mm]
  \setcounter{solstep}{0}
  \begin{solutionblock}
    \textbf{答案：}(1) 圆心距为 $6$，位置关系为**外离**；(2) $t = 3$ 或 $5$ 或 $7$ 或 $9$。\par\medskip
    \solstep{求初始圆心距和关系}{圆心距 $d = |3 - (-3)| = 6$。因为 $6 > 1+2=3$，所以两圆初始时**外离**。}
    \solstep{表示平移后的圆心距}{平移 $t$ 秒后，新圆心为 $A'(t-3, 0)$，圆心距 $d(t) = |3 - (t-3)| = |6-t|$。}
    \solstep{外切情况求解}{外切时，$|6-t| = 1+2=3 \implies 6-t = \pm 3 \implies t = 3$ 或 $t = 9$。}
    \solstep{内切情况求解}{内切时，$|6-t| = 2-1=1 \implies 6-t = \pm 1 \implies t = 5$ 或 $t = 7$。}
  \end{solutionblock}
\end{problem}


\needspace{8\baselineskip}

\begin{problem}[points=10]
  已知 $\odot O_1$ 与 $\odot O_2$ 相切，圆心距 $d = 5$。若这两圆的半径 $R$ 和 $r$（$R \ge r$）是关于 $x$ 的一元二次方程 $x^2 - mx + 6 = 0$（$m > 0$）的两个实数根，求 $m$ 的值。

  \answerarea[70mm]
  \setcounter{solstep}{0}
  \begin{solutionblock}
    \textbf{答案：}$m = 5$ 或 $7$\par\medskip
    \solstep{应用韦达定理}{$R+r = m$，$R \cdot r = 6$。判别式 $\Delta = m^2 - 24 \ge 0$。}
    \solstep{外切讨论}{外切时 $R+r = 5 \implies m = 5$。此时 $\Delta = 25-24=1 \ge 0$，有效。}
    \solstep{内切讨论}{内切时 $|R-r| = 5 \implies (R-r)^2 = 25$。代入恒等式 $(R+r)^2 - 4Rr = 25 \implies m^2 - 24 = 25 \implies m^2 = 49 \implies m = 7$。此时 $\Delta = 49-24=25 \ge 0$，有效。}
  \end{solutionblock}
\end{problem}


\needspace{8\baselineskip}

\begin{problem}[points=10]
  如图，在边长为 $4$ 的正方形 $ABCD$ 中，点 $E$ 是边 $BC$ 上的一点。以点 $E$ 为圆心、$EC$ 为半径的半圆 $\odot E$ 与以点 $A$ 为圆心、$AB$ 为半径的圆弧 $\odot A$ 外切。
\begin{enumerate}[label=(\arabic*)]
  \item 联结 $AE$，求证：$AE = 4 + EC$；
  \item 求线段 $EC$ 的长度；
  \item 求 $\sin\angle EAB$ 的值。
\end{enumerate}

  \answerarea[80mm]
  \setcounter{solstep}{0}
  \begin{solutionblock}
    \textbf{答案：}(1) 证明见解析；(2) $EC = 1$；(3) $\sin\angle EAB = \dfrac{3}{5}$。\par\medskip
    \solstep{第(1)问证明}{因为两圆外切，连心线段等于半径之和，即 $AE = R_A + R_E$。因为 $\odot A$ 过 $B$，半径 $R_A = AB = 4$；$\odot E$ 半径为 $EC$。所以 $AE = 4 + EC$。}
    \solstep{第(2)问列方程}{设 $EC = x$，则 $AE = 4+x$，$BE = 4-x$。在 $\mathrm{Rt}\triangle ABE$ 中，由 $AB^2 + BE^2 = AE^2$ 得 $16 + (4-x)^2 = (4+x)^2$。}
    \solstep{求解 x}{化简方程：$16 + 16 - 8x + x^2 = 16 + 8x + x^2 \implies 16 = 16x \implies x = 1$。所以 $EC = 1$。}
    \solstep{第(3)问求正弦值}{由 $x = 1$ 得 $BE = 3$，$AE = 5$。$\sin\angle EAB = \dfrac{BE}{AE} = \dfrac{3}{5}$。}
  \end{solutionblock}
\end{problem}


\clearpage
\section*{参考答案}


\begin{keyideabox}
  \textbf{选择题}
  1. C
2. B

\end{keyideabox}



\begin{keyideabox}
  \textbf{填空题}
  1. $4$ 或 $10$
2. $2$ 或 $8$
3. $d > 8$ 或 $0 \le d < 2$

\end{keyideabox}



\begin{keyideabox}
  \textbf{解答题}
  第 1 题：
(1) 圆心距为 $6$，位置关系为外离；
(2) $t = 3$ 或 $5$ 或 $7$ 或 $9$ 秒。

第 2 题：
$m = 5$ 或 $7$。

第 3 题：
(1) 证明：因为外切，$AE = R_A + R_E = 4 + EC$；
(2) $EC = 1$；
(3) $\sin\angle EAB = \dfrac{3}{5}$。

\end{keyideabox}



\end{document}
